% ECONOMICS_PROBLEM: {"id": "C73-19", "title": "Polynomial algorithms for games with ergodic payoff", "jel_code": "C73", "source_id": "OP-0164"}
% FIRST_POSTED: 2026-10-09
% ATTEMPT_STATUS: unresolved
% ENTRY_KIND: research_attempt
% !TeX program = pdflatex
% Self-contained source: pdflatex twice; no fontspec or external bibliography.
\documentclass[11pt,a4paper]{article}
\usepackage[margin=24mm,headheight=14pt]{geometry}
\usepackage[T1]{fontenc}
\usepackage[utf8]{inputenc}
\usepackage{lmodern,microtype}
\usepackage{amsmath,amssymb,mathtools,amsthm}
\usepackage{enumitem,booktabs,longtable,array,xurl,listings,needspace}
\usepackage[hidelinks,unicode]{hyperref}
\hypersetup{pdftitle={Economic Theory: Proof Attempts and Adversarial Checks},pdfauthor={Research notes},pdfsubject={Eleven user-supplied mathematical problems}}
\newcommand{\NE}{\operatorname{NE}}
\DeclareMathOperator*{\argmax}{arg\,max}
\DeclareMathOperator*{\argmin}{arg\,min}
\newtheorem{theorem}{Theorem}[section]
\newtheorem{lemma}[theorem]{Lemma}
\newtheorem{proposition}[theorem]{Proposition}
\newtheorem{corollary}[theorem]{Corollary}
\theoremstyle{definition}
\newtheorem{definition}[theorem]{Definition}
\theoremstyle{remark}
\newtheorem*{remark}{Remark}
\setlength{\parindent}{0pt}
\setlength{\parskip}{5pt plus 1pt minus 1pt}
\setlength{\emergencystretch}{2em}
\setcounter{tocdepth}{1}
\setlist{leftmargin=*,itemsep=2pt,topsep=4pt}
\lstset{basicstyle=\small\ttfamily,columns=fullflexible,keepspaces=true,
 breaklines=true,breakatwhitespace=false,showstringspaces=false,
 frame=single,framerule=0.25pt,aboveskip=8pt,belowskip=8pt,xleftmargin=3pt,xrightmargin=3pt}
\allowdisplaybreaks[1]
\raggedbottom

\begin{document}
{\large\bfseries Research attempt and adversarial verification}\par
9 October 2026\hfill Original-target status: \textbf{UNRESOLVED}\par
\section{C73-19: Polynomial algorithms under state-independent mean payoff}
\label{sec:C73-19}
\noindent\textbf{Input:} \texttt{C73-19.tex}. \quad\textbf{Tools:} web browsing [YES]; Python computation [YES].
\subsection*{1) FORMAL RESTATEMENT}
An input lists a nonempty finite state set $S$, a controlling player (Max or Min) at every state, nonempty finite action sets $A(s)$, rational rewards $r(s,a)$, rational stochastic kernels $P(\cdot\mid s,a)$, and a rational threshold $q$, all explicitly in binary. Single-action states allow uncontrolled randomness. For arbitrary behavioral strategies, time starts at zero and
\[
 g_s(\sigma,\tau)=\mathbb E_{s,\sigma,\tau}
 \left[\liminf_{T\to\infty}\frac1T\sum_{t=0}^{T-1}r(s_t,a_t)\right].
\]
The promise is $\exists\rho\in\mathbb R\ \forall s\in S$, $\sup_\sigma\inf_\tau g_s=\inf_\tau\sup_\sigma g_s=\rho$. The question is whether there exist a deterministic algorithm and a polynomial $p$ such that every promised input of bit length $L$ is correctly classified according to $\rho\geq q$ within $p(L)$ steps. Behavior on unpromised inputs need not be correct.

No correction is needed. In particular, the promise does not assert irreducibility, a mixing-time bound, or equality of different orders of expectation and $\liminf$. A counterexample to a proposed algorithm is not a disproof of polynomial-time solvability.

\subsection*{2) QUICK LITERATURE/CONTEXT CHECK}
The TROPICAL activity report \cite{tropical} lists polynomial algorithms for games with ergodic payment as a research direction; its brief statement does not fix the exact promise in this file. Hansen and Nie \cite{hansen-nie} discuss the unresolved zero-sum discounted/simple-stochastic-game complexity question. The reduction below is proved directly for the displayed state-independent-value promise and is not based on interpreting all possible meanings of ``ergodic'' as equivalent. The check is dated 9 October 2026.

\subsection*{3) ATTACK PLAN}
\textbf{Type and tools.} This is exact bit complexity for stochastic games. Relevant tools are: Bellman contraction (discounted values); additive bias certificates (mean payoff); bounded-difference martingales (pathwise averages); rational reductions (preserve input size); regenerative resets (force a common value); elementary hitting-time bounds (stress-test iteration); stationary policies (finite certificates).

\textbf{Proof tracks.} Seek a polynomial bias computation or reduce the promised class to a known tractable subclass. \textbf{Disproof tracks.} Embed a discounted game using resets; force very slow convergence with a tiny binary probability; check whether a scalar Bellman eigenvalue actually controls the stipulated pathwise payoff. The reset reduction and a complete pathwise verification lemma are the strongest achieved route.

\subsection*{4) WORK}
\begin{lemma}[Pathwise additive-bias certificate]
Suppose $h:S\to\mathbb R$ and $\rho\in\mathbb R$ satisfy
\[
 \rho+h(s)=
 \begin{cases}
 \max_{a\in A(s)}[r(s,a)+\sum_{s'}P(s'\mid s,a)h(s')],&s\text{ Max},\\
 \min_{a\in A(s)}[r(s,a)+\sum_{s'}P(s'\mid s,a)h(s')],&s\text{ Min}.
 \end{cases}
\]
Choosing an attaining action at each controlled state gives stationary strategies that guarantee the common value $\rho$, from every initial state, for the expectation of the pathwise $\liminf$ criterion in the input.
\end{lemma}
\begin{proof}
Fix Max's maximizing selector and any behavioral Min strategy. Conditional on the past and the chosen action, every step satisfies
\[
 r(s_t,a_t)+\mathbb E[h(s_{t+1})\mid\mathcal F_t,a_t]-h(s_t)\geq\rho.
\]
Here $\mathcal F_t$ contains all information before the transition; using a filtration including the current action removes any issue with randomized actions. Let
$D_{t+1}=h(s_{t+1})-\mathbb E[h(s_{t+1})\mid\mathcal F_t,a_t]$ and $M_T=\sum_{t<T}D_{t+1}$. These are bounded martingale differences, with $|D_{t+1}|\leq2\|h\|_\infty$. Summing gives
\[
 \sum_{t<T}(r(s_t,a_t)-\rho)
 \geq h(s_0)-h(s_T)+M_T.
\]
For every $\eta>0$, the bounded-difference martingale inequality gives a bound $2\exp[-\eta^2T/(8\|h\|_\infty^2)]$ on $\Pr(|M_T|>\eta T)$ when $\|h\|_\infty>0$. The probabilities are summable; the elementary Borel--Cantelli lemma implies $M_T/T\to0$ almost surely. If $h=0$ this conclusion is immediate. The bounded telescoping term divided by $T$ also tends to zero. Thus the pathwise $\liminf$ is at least $\rho$ almost surely, and its expectation is at least $\rho$ because rewards are bounded.

Fixing Min's minimizing selector reverses every inequality for arbitrary Max play and gives a pathwise $\limsup$ at most $\rho$. Therefore its pathwise $\liminf$ and its expectation are at most $\rho$. The general inequality maxmin $\leq$ minimax and these two guarantees prove equality. No exchange of expectation and $\liminf$ was used.
\end{proof}

\begin{proposition}[Discounted games embed in the promised class]
Given a finite rational turn-based zero-sum discounted game with $0<\beta<1$, rewards $r$, transitions $P$, and initial state $s_*$, form
\[
 P'(s'\mid s,a)=\beta P(s'\mid s,a)+(1-\beta)\mathbf1_{\{s'=s_*\}}.
\]
Keep the rewards and controlling players unchanged. The new undiscounted game satisfies the input's promise and its common mean value is the normalized discounted value from $s_*$ in the original game. This construction has polynomial bit complexity.
\end{proposition}
\begin{proof}
The discounted Bellman map $T_\beta$ takes a vector $z$ to the maximum or minimum, according to the state's owner, of $r(s,a)+\beta Pz(s,a)$. A maximum or minimum of finitely many real numbers changes by at most their largest coordinatewise change. Hence $\|T_\beta z-T_\beta w\|_\infty\leq\beta\|z-w\|_\infty$. Iteration is Cauchy by the geometric bound on successive increments, and completeness of $\mathbb R^S$ supplies a unique fixed point $h$.

The fixed point is the unnormalized discounted value. To verify this rather than assume it, choose a maximizing action in its equation at each Max state. Against arbitrary Min actions,
$r(s_t,a_t)+\beta\mathbb E h(s_{t+1})\geq h(s_t)$ conditionally. Multiply by $\beta^t$, sum expectations to time $T-1$, and use the bound $\beta^T\|h\|_\infty\to0$. This guarantees expected discounted sum at least $h(s_0)$. Min's selector gives the reverse inequality. The normalized criterion is consequently $(1-\beta)h(s_0)$.

For the reset game,
\[
 r+P'h=r+\beta Ph+(1-\beta)h(s_*).
\]
The last term is the same for every state and action. Taking the appropriate maximum or minimum yields
\[
 \max/\min\,(r+P'h)=h(s)+\rho,
 \qquad \rho=(1-\beta)h(s_*).
\]
The preceding lemma proves the mean-payoff promise and value from every initial state. Each output probability is a sum of at most two products of rational input numbers. Its numerator and denominator have bit length at most a constant multiple of the corresponding input lengths, and the number of states and actions is unchanged. Normalized discounted threshold $q$ maps to the same threshold $q$. For an unnormalized threshold, multiply it by $1-\beta$. Thus the reduction is polynomial.
\end{proof}

\paragraph{What this reduction proves.}
A polynomial exact algorithm for the promised class would give one for arbitrary rational discounted turn-based zero-sum games. It does not prove that either problem lacks such an algorithm. Resetting to $s_*$ does not imply every state is reachable from $s_*$; the proposition claims state-independent value, exactly the specified promise, not a stronger irreducibility property.

\begin{proposition}[The promise does not give polynomial horizon convergence]
For each integer $b\geq1$, there is a two-state uncontrolled promised game with $O(b)$-bit data, common value one, and expected average reward at horizon $T$ at most $2^{-b}(T-1)/2$ when started in its first state.
\end{proposition}
\begin{proof}
At state $a$ the reward is zero; transition to absorbing $b_1$ has probability $\delta=2^{-b}$ and otherwise the state remains $a$. At $b_1$ the reward is one. The hitting time is finite almost surely because the probability of remaining in $a$ for $t$ transitions is $(1-\delta)^t\to0$. Every sample path that hits $b_1$ has average reward tending to one; boundedness gives the required common expected pathwise value. Starting at $a$, $\Pr(s_t=b_1)=1-(1-\delta)^t\leq t\delta$, by a union bound on the first $t$ transitions. Summing over $0\leq t<T$ and dividing by $T$ gives the assertion. For $T\leq2^{b-1}$ the expected average is less than $1/4$. Thus even constant-accuracy finite-horizon substitution can require exponentially many stages in the bit length on this family.
\end{proof}

\paragraph{Tiny cases and exact calculation.}
One state reduces to the owner's maximum or minimum immediate reward. For a nontrivial reset test use $\beta=2/3$, states $0$ (Max) and $1$ (Min), and actions listed as (reward, next state):
\[
 A(0)=\{(0,0),(1,1)\},\qquad A(1)=\{(0,0),(2,1)\}.
\]
The discounted fixed point is $h=(9/5,6/5)$, since the two action values are $(6/5,9/5)$ at state zero and $(6/5,14/5)$ at state one. Reset to zero gives $\rho=3/5$. The reset action expressions $r+P'h$ are $(9/5,12/5)$ and $(9/5,17/5)$; their maximum and minimum equal $h+\rho$. The script checks these fractions exactly.
\begin{lstlisting}
for state in states:
    rhs = [reward[state,a] + sum(Preset[state,a,z]*h[z]
                               for z in states)
           for a in actions[state]]
    assert owner_extremum(rhs) == h[state] + rho
\end{lstlisting}

\subsection*{5) VERIFICATION}
The bias proof covers arbitrary behavioral deviations, not just stationary competitors. Bounded martingale increments hold because the state set and $h$ are finite, even when the input bit lengths are large. The reset construction has no vanishing-probability limit: it is an exact finite transformation at the input discount factor. The slow-convergence example meets the promise despite a transient state. Neither a poor finite-horizon algorithm nor a complexity-preserving embedding is an unconditional complexity lower bound. The threshold's binary length remains part of every running-time assertion.

\Needspace{8\baselineskip}
\subsection*{6) FINAL}
\textbf{LABEL: UNRESOLVED} for polynomial-time exact threshold decision on every promised input.

\textbf{Strongest checked partial result.} An exact polynomial reduction from rational discounted turn-based zero-sum games to the stated promise, supported by a pathwise bias-certificate proof. The two-state family also rigorously rules out a proposed proof based solely on uniformly polynomial finite-horizon convergence.

\textbf{Exact first gap.} There is no proved algorithm here that computes or separates the additive eigenvalue in time polynomial in input bit length on the reset subclass, let alone on every promised game.

\textbf{Top three next moves.} (1) Seek a strongly bounded number of policy changes for reset Bellman operators, independent of $(1-\beta)^{-1}$. (2) Construct exact rational primal/dual certificates that can be found without enumerating policy pairs. (3) Determine whether the scalar-value promise permits a polynomial decomposition when a bounded additive bias does not exist.

\textbf{Minimal counterexample perspective.} Complexity cannot be refuted by one finite game. A negative result would need a family and a justified computational lower-bound framework. The two-state rare-transition family is only a counterexample to a horizon-based algorithmic shortcut.

\paragraph{Reproducibility and scope.} This is one of eleven companion reports. The bundle includes \texttt{verify.py} and its executed results. Finite experiments supplement the proofs; they are not proofs of asymptotic claims. Literature coverage is targeted, and no novelty or formal proof-assistant certification is claimed.

\begin{thebibliography}{99}
\small
\bibitem{tropical}
Inria, TROPICAL research team. \emph{TROPICAL Activity Report 2024}.
Section 3.1, printed p.~4 (PDF page 6).
\url{https://radar.inria.fr/rapportsactivite/RA2024/tropical/TROPICAL-RA-2024.pdf}.

\bibitem{hansen-nie}
Kristoffer Arnsfelt Hansen and Xinhao Nie.
\emph{On the Complexity of Stationary Nash Equilibria in Discounted Perfect Information Stochastic Games}.
arXiv:2510.11550, version 1 (2025).
\url{https://arxiv.org/html/2510.11550v1}.

\end{thebibliography}

\end{document}
